\section{スケジュールと担当}

図\ref{suke}は計画段階でのスケジュールである.
まず音のシステムはマイクが送信した音を送信した通りに受け取らないものが出てきてしまったため,16進数から音階を減らして正しく受け取ることのできる音階を使用した4進数に変更したためプログラムの作成が遅れてしまった.
またマイクが壊れてしまったので取り付け直し,正しく受け取らない音階の原因を探るためにスピーカーも取り付け直した.
次に光のシステムはすべて計画通りに行えた.
最後に球のシステムは送信側のモータは特に問題点なく進められたため,計画より早く取り組めた.
しかし受信側の照度センサが上手く反応しなかったため,プログラムの調整と装置の調整を9,10週目でも取り組んだ.
結果,実際のスケジュールは図\ref{suke_2}のようになった.

\begin{figure}[htbp]
\begin{center}
\includegraphics[width=\textwidth,clip]{image/sukezyuru.png}
\caption{計画段階でのスケジュールと担当}
\label{suke}
\end{center}
\end{figure}

\begin{figure}[htbp]
\begin{center}
\includegraphics[width=\textwidth,clip]{image/sukezyuru_2.png}
\caption{実際のスケジュールと担当}
\label{suke_2}
\end{center}
\end{figure}